\documentclass[10pt,a4paper]{amsart}
\usepackage{amsmath,amssymb,amsthm,graphicx}
% Retained minimal preamble for the Stacks Project chapter flat.tex (More on Flatness).
% Verbatim excerpt of the chapter preamble: Xy-pic setup (native preamble lines 14-18) and the native
% theorem declarations with the mathematical macros (native preamble lines 154-217). Copied mechanically;
% nothing added, removed or redefined. The wrapper class and the amsmath/amssymb/amsthm/graphicx base
% packages are supplied by the assembler, so the class/externaldocument part of the native preamble is omitted.
% For commutative diagrams we use Xy-pic
\usepackage[all]{xy}

% We use 2cell for 2-commutative diagrams.
\xyoption{2cell}
% Theorem environments.
%
\theoremstyle{plain}
\newtheorem{theorem}[subsection]{Theorem}
\newtheorem{proposition}[subsection]{Proposition}
\newtheorem{lemma}[subsection]{Lemma}

\theoremstyle{definition}
\newtheorem{definition}[subsection]{Definition}
\newtheorem{example}[subsection]{Example}
\newtheorem{exercise}[subsection]{Exercise}
\newtheorem{situation}[subsection]{Situation}

\theoremstyle{remark}
\newtheorem{remark}[subsection]{Remark}
\newtheorem{remarks}[subsection]{Remarks}

\numberwithin{equation}{subsection}

% Macros
%
\def\lim{\mathop{\mathrm{lim}}\nolimits}
\def\colim{\mathop{\mathrm{colim}}\nolimits}
\def\Spec{\mathop{\mathrm{Spec}}}
\def\Hom{\mathop{\mathrm{Hom}}\nolimits}
\def\Ext{\mathop{\mathrm{Ext}}\nolimits}
\def\SheafHom{\mathop{\mathcal{H}\!\mathit{om}}\nolimits}
\def\SheafExt{\mathop{\mathcal{E}\!\mathit{xt}}\nolimits}
\def\Sch{\mathit{Sch}}
\def\Mor{\mathop{\mathrm{Mor}}\nolimits}
\def\Ob{\mathop{\mathrm{Ob}}\nolimits}
\def\Sh{\mathop{\mathit{Sh}}\nolimits}
\def\NL{\mathop{N\!L}\nolimits}
\def\CH{\mathop{\mathrm{CH}}\nolimits}
\def\proetale{{pro\text{-}\acute{e}tale}}
\def\etale{{\acute{e}tale}}
\def\QCoh{\mathit{QCoh}}
\def\Ker{\mathop{\mathrm{Ker}}}
\def\Im{\mathop{\mathrm{Im}}}
\def\Coker{\mathop{\mathrm{Coker}}}
\def\Coim{\mathop{\mathrm{Coim}}}

% Boxtimes
%
\DeclareMathSymbol{\boxtimes}{\mathbin}{AMSa}{"02}

%
% Macros for moduli stacks/spaces
%
\def\QCohstack{\mathcal{QC}\!\mathit{oh}}
\def\Cohstack{\mathcal{C}\!\mathit{oh}}
\def\Spacesstack{\mathcal{S}\!\mathit{paces}}
\def\Quotfunctor{\mathrm{Quot}}
\def\Hilbfunctor{\mathrm{Hilb}}
\def\Curvesstack{\mathcal{C}\!\mathit{urves}}
\def\Polarizedstack{\mathcal{P}\!\mathit{olarized}}
\def\Complexesstack{\mathcal{C}\!\mathit{omplexes}}
% \Pic is the operator that assigns to X its picard group, usage \Pic(X)
% \Picardstack_{X/B} denotes the Picard stack of X over B
% \Picardfunctor_{X/B} denotes the Picard functor of X over B
\def\Pic{\mathop{\mathrm{Pic}}\nolimits}
\def\Picardstack{\mathcal{P}\!\mathit{ic}}
\def\Picardfunctor{\mathrm{Pic}}
\def\Deformationcategory{\mathcal{D}\!\mathit{ef}}

% --- appended 2026-10-12 (Chief): the project own definitions of the reference/slogan/history
% display environments, copied mechanically from the shared Stacks preamble lines 8-12
% (\usepackage{verbatim} plus the three \newenvironment lines, which the chapter preamble keeps
% verbatim from the project preamble). Nothing added, removed or redefined: these environments
% carry no mathematical content and the project itself typesets them as comments. Without these
% three lines a counted statement whose own text uses them fails with Environment slogan undefined.
\usepackage{verbatim}
\newenvironment{reference}{\comment}{\endcomment}
\newenvironment{slogan}{\comment}{\endcomment}
\newenvironment{history}{\comment}{\endcomment}
\usepackage{cleveref}
\crefname{originallemma}{original lemma}{original lemmas}
\Crefname{originallemma}{Original lemma}{Original lemmas}
\makeatletter
\newlabel{algebra-lemma-extension}{{0519}{source}{}{sourcepointer-1}{}}
\newlabel{algebra-lemma-extension@cref}{{[originallemma][0][]0519}{[1][0][]source}{}{}{}}
\makeatother
\crefname{originallemma}{original lemma}{original lemmas}
\Crefname{originallemma}{Original lemma}{Original lemmas}
\makeatletter
\newlabel{algebra-lemma-finite-flat-local}{{00NZ}{source}{}{sourcepointer-2}{}}
\newlabel{algebra-lemma-finite-flat-local@cref}{{[originallemma][0][]00NZ}{[1][0][]source}{}{}{}}
\makeatother
\crefname{originallemma}{original lemma}{original lemmas}
\Crefname{originallemma}{Original lemma}{Original lemmas}
\makeatletter
\newlabel{algebra-lemma-flat-ses}{{00HM}{source}{}{sourcepointer-3}{}}
\newlabel{algebra-lemma-flat-ses@cref}{{[originallemma][0][]00HM}{[1][0][]source}{}{}{}}
\makeatother
\crefname{originallemma}{original lemma}{original lemmas}
\Crefname{originallemma}{Original lemma}{Original lemmas}
\makeatletter
\newlabel{descent-lemma-universally-injective-etale-open-immersion}{{02LC}{source}{}{sourcepointer-4}{}}
\newlabel{descent-lemma-universally-injective-etale-open-immersion@cref}{{[originallemma][0][]02LC}{[1][0][]source}{}{}{}}
\makeatother
\crefname{originallemma}{original lemma}{original lemmas}
\Crefname{originallemma}{Original lemma}{Original lemmas}
\makeatletter
\newlabel{lemma-devissage-finite-presentation}{{057T}{source}{}{sourcepointer-5}{}}
\newlabel{lemma-devissage-finite-presentation@cref}{{[originallemma][0][]057T}{[1][0][]source}{}{}{}}
\makeatother
\crefname{originallemma}{original lemma}{original lemmas}
\Crefname{originallemma}{Original lemma}{Original lemmas}
\makeatletter
\newlabel{lemma-devissage-flat}{{057U}{source}{}{sourcepointer-6}{}}
\newlabel{lemma-devissage-flat@cref}{{[originallemma][0][]057U}{[1][0][]source}{}{}{}}
\makeatother
\crefname{originallemma}{original lemma}{original lemmas}
\Crefname{originallemma}{Original lemma}{Original lemmas}
\makeatletter
\newlabel{lemma-elementary-devissage}{{057S}{source}{}{sourcepointer-7}{}}
\newlabel{lemma-elementary-devissage@cref}{{[originallemma][0][]057S}{[1][0][]source}{}{}{}}
\makeatother
\crefname{originallemma}{original lemma}{original lemmas}
\Crefname{originallemma}{Original lemma}{Original lemmas}
\makeatletter
\newlabel{lemma-etale-at-point}{{05H2}{source}{}{sourcepointer-8}{}}
\newlabel{lemma-etale-at-point@cref}{{[originallemma][0][]05H2}{[1][0][]source}{}{}{}}
\makeatother
\crefname{originallemma}{original lemma}{original lemmas}
\Crefname{originallemma}{Original lemma}{Original lemmas}
\makeatletter
\newlabel{lemma-induction-step}{{05I3}{source}{}{sourcepointer-9}{}}
\newlabel{lemma-induction-step@cref}{{[originallemma][0][]05I3}{[1][0][]source}{}{}{}}
\makeatother
\title{Local structure of a finite type quasi-coherent sheaf flat at a point, by elementary etale neighbourhood and open subscheme}
\author{The Stacks Project Authors}
\date{stacks-project:flat.tex:local-snapshot-20261011; GNU FDL 1.2 or later}
\begin{document}
\maketitle
\noindent\textit{Editorial note.} All delivered mathematical text is the author's own, copied from the pinned Stacks Project chapter file flat.tex; no statement or proof step was rewritten, reordered or removed.\par\medskip
\noindent\textit{Editorial note.} Preamble provenance: the chapter's own retained minimal preamble (sources/stacks-flat/retained-preamble.tex, lines 1-85) is a verbatim excerpt of the native Xy-pic setup and of the native theorem declarations and mathematical macros, followed by the project's own reference/slogan/history display-environment definitions copied mechanically from the shared preamble; the wrapper class and the amsmath/amssymb/amsthm/graphicx base packages are supplied by the assembler, and extra\_packages is left empty because everything needed is already in the retained preamble. No macro was invented or redefined.\par\medskip
\noindent\textit{Editorial note.} Same-chapter pointers are located at their \textbackslash{}label lines inside the pinned chapter file; cross-chapter pointers are located in the project's own tag list staged in the snapshot (sources/stacks-chow/tags.txt), at the verbatim lines carrying their canonical tags.\par\medskip
\noindent\textit{Editorial note.} The source gives the proposition two proofs, the first beginning at source line 2498 under the heading 'First proof' and saying that it is longer but does not use the existence of a complete devissage. The counted proof span is that first proof, which is a complete proof of the statement; the second proof immediately following was not taken, so that the counted pair stays one statement with one complete proof. The author's remark that the problem is local around the two points is used verbatim as the first reduction step.\par\medskip
\noindent\textit{Allowed standard prerequisites (editorial).}\begin{itemize}
\item Elementary devissage of a finite type quasi-coherent sheaf relative to a morphism. Same chapter flat.tex, label lemma-elementary-devissage, invoked at proof lines 2528 and 2550. Produces the elementary etale neighbourhoods, the diagram with Z' -> Y' -> S', the point z' and the finite type module G which is the starting point of the reduction.
\item Descent of finite presentation, flatness and etaleness along the devissage. Same chapter flat.tex, labels lemma-devissage-finite-presentation (proof line 2552), lemma-devissage-flat (proof line 2553) and lemma-etale-at-point (proof line 2544). Shows the result for Y' -> S' with the pushed forward module implies the result for X -> S with F, and keeps the base change map etale after replacing S by a localisation.
\item The induction step producing a finite free presentation with flat cokernel, and flatness and finite presentation in a short exact sequence. Same chapter flat.tex, label lemma-induction-step, invoked at proof line 2580; Algebra, labels algebra-lemma-finite-flat-local (proof line 2569), algebra-lemma-flat-ses (proof line 2611) and algebra-lemma-extension (proof line 2614); Descent, label descent-lemma-universally-injective-etale-open-immersion (proof line 2564). Gives the finite free module in the base case, supplies the map alpha with isomorphic reduction and flat cokernel in the induction step, and lets flatness and finite presentation be inferred for the middle term of the resulting short exact sequence.
\end{itemize}\medskip
\ifdefined\hypertarget\hypertarget{sourcepointer-1}{}\fi
\noindent\textit{Source reference.} Original lemma 0519: sources/stacks-chow/tags.txt, source lines 6527--6527. Navigation pointer for \textbackslash{}ref\{algebra-lemma-extension\} used inside the counted span (line 2614); the argument that consumes it is retained in the span.\par\medskip
\ifdefined\hypertarget\hypertarget{sourcepointer-2}{}\fi
\noindent\textit{Source reference.} Original lemma 00NZ: sources/stacks-chow/tags.txt, source lines 877--877. Navigation pointer for \textbackslash{}ref\{algebra-lemma-finite-flat-local\} used inside the counted span (line 2569); the argument that consumes it is retained in the span.\par\medskip
\ifdefined\hypertarget\hypertarget{sourcepointer-3}{}\fi
\noindent\textit{Source reference.} Original lemma 00HM: sources/stacks-chow/tags.txt, source lines 648--648. Navigation pointer for \textbackslash{}ref\{algebra-lemma-flat-ses\} used inside the counted span (line 2611); the argument that consumes it is retained in the span.\par\medskip
\ifdefined\hypertarget\hypertarget{sourcepointer-4}{}\fi
\noindent\textit{Source reference.} Original lemma 02LC: sources/stacks-chow/tags.txt, source lines 3403--3403. Navigation pointer for \textbackslash{}ref\{descent-lemma-universally-injective-etale-open-immersion\} used inside the counted span (line 2564); the argument that consumes it is retained in the span.\par\medskip
\ifdefined\hypertarget\hypertarget{sourcepointer-5}{}\fi
\noindent\textit{Source reference.} Original lemma 057T: sources/stacks-flat/pinned-original-source.tex, source lines 598--598. Navigation pointer for \textbackslash{}ref\{lemma-devissage-finite-presentation\} used inside the counted span (line 2552); the argument that consumes it is retained in the span.\par\medskip
\ifdefined\hypertarget\hypertarget{sourcepointer-6}{}\fi
\noindent\textit{Source reference.} Original lemma 057U: sources/stacks-flat/pinned-original-source.tex, source lines 688--688. Navigation pointer for \textbackslash{}ref\{lemma-devissage-flat\} used inside the counted span (line 2553); the argument that consumes it is retained in the span.\par\medskip
\ifdefined\hypertarget\hypertarget{sourcepointer-7}{}\fi
\noindent\textit{Source reference.} Original lemma 057S: sources/stacks-flat/pinned-original-source.tex, source lines 494--494. Navigation pointer for \textbackslash{}ref\{lemma-elementary-devissage\} used inside the counted span (line 2528); the argument that consumes it is retained in the span.\par\medskip
\ifdefined\hypertarget\hypertarget{sourcepointer-8}{}\fi
\noindent\textit{Source reference.} Original lemma 05H2: sources/stacks-flat/pinned-original-source.tex, source lines 97--97. Navigation pointer for \textbackslash{}ref\{lemma-etale-at-point\} used inside the counted span (line 2544); the argument that consumes it is retained in the span.\par\medskip
\ifdefined\hypertarget\hypertarget{sourcepointer-9}{}\fi
\noindent\textit{Source reference.} Original lemma 05I3: sources/stacks-flat/pinned-original-source.tex, source lines 2356--2356. Navigation pointer for \textbackslash{}ref\{lemma-induction-step\} used inside the counted span (line 2580); the argument that consumes it is retained in the span.\par\medskip
\begin{proposition}
\label{proposition-finite-type-flat-at-point}
Let $f : X \to S$ be a morphism of schemes.
Let $\mathcal{F}$ be a quasi-coherent sheaf on $X$.
Let $x \in X$ with image $s \in S$.
Assume that
\begin{enumerate}
\item $f$ is locally of finite presentation,
\item $\mathcal{F}$ is of finite type, and
\item $\mathcal{F}$ is flat at $x$ over $S$.
\end{enumerate}
Then there exists an elementary \'etale neighbourhood $(S', s') \to (S, s)$
and an open subscheme
$$
V \subset X \times_S \Spec(\mathcal{O}_{S', s'})
$$
which contains the unique point of
$X \times_S \Spec(\mathcal{O}_{S', s'})$ mapping to $x$
such that the pullback of $\mathcal{F}$ to $V$ is an $\mathcal{O}_V$-module
of finite presentation and flat over $\mathcal{O}_{S', s'}$.
\end{proposition}

\begin{proof}[First proof]
This proof is longer but does not use the existence of a complete d\'evissage.
The problem is local around $x$ and $s$, hence we may assume that $X$
and $S$ are affine. During the proof we will finitely many times replace
$S$ by an elementary \'etale neighbourhood of $(S, s)$. The goal is then to find
(after such a replacement) an open
$V \subset X \times_S \Spec(\mathcal{O}_{S, s})$ containing $x$
such that $\mathcal{F}|_V$ is flat over $S$ and finitely presented.
Of course we may also replace $S$ by $\Spec(\mathcal{O}_{S, s})$
at any point of the proof, i.e., we may assume $S$ is a local scheme.
We will prove the proposition by induction on the integer
$n = \dim_x(\text{Supp}(\mathcal{F}_s))$.

\medskip\noindent
We can choose
\begin{enumerate}
\item elementary \'etale neighbourhoods $g : (X', x') \to (X, x)$,
$e : (S', s') \to (S, s)$,
\item a commutative diagram
$$
\xymatrix{
X \ar[dd]_f & X' \ar[dd] \ar[l]^g & Z' \ar[l]^i \ar[d]^\pi \\
& & Y' \ar[d]^h \\
S & S' \ar[l]_e & S' \ar@{=}[l]
}
$$
\item a point $z' \in Z'$ with $i(z') = x'$, $y' = \pi(z')$, $h(y') = s'$,
\item a finite type quasi-coherent $\mathcal{O}_{Z'}$-module $\mathcal{G}$,
\end{enumerate}
as in
Lemma \ref{lemma-elementary-devissage}.
We are going to replace $S$ by $\Spec(\mathcal{O}_{S', s'})$, see
remarks in first paragraph of the proof. Consider the diagram
$$
\xymatrix{
X_{\mathcal{O}_{S', s'}} \ar[ddr]_f &
X'_{\mathcal{O}_{S', s'}} \ar[dd] \ar[l]^g &
Z'_{\mathcal{O}_{S', s'}} \ar[l]^i \ar[d]^\pi \\
& & Y'_{\mathcal{O}_{S', s'}} \ar[dl]^h \\
& \Spec(\mathcal{O}_{S', s'})
}
$$
Here we have base changed the schemes $X', Z', Y'$ over $S'$ via
$\Spec(\mathcal{O}_{S', s'}) \to S'$ and the scheme $X$ over $S$ via
$\Spec(\mathcal{O}_{S', s'}) \to S$. It is still the case that
$g$ is \'etale, see
Lemma \ref{lemma-etale-at-point}.
After replacing $X$ by $X_{\mathcal{O}_{S', s'}}$,
$X'$ by $X'_{\mathcal{O}_{S', s'}}$,
$Z'$ by $Z'_{\mathcal{O}_{S', s'}}$, and
$Y'$ by $Y'_{\mathcal{O}_{S', s'}}$
we may assume we have a diagram as
Lemma \ref{lemma-elementary-devissage}
where in addition $S = S'$ is a local scheme with closed point $s$. By
Lemmas \ref{lemma-devissage-finite-presentation} and
\ref{lemma-devissage-flat}
the result for $Y' \to S$, the sheaf $\pi_*\mathcal{G}$, and the
point $y'$ implies the result for $X \to S$, $\mathcal{F}$ and $x$.
Hence we may assume that $S$ is local and $X \to S$ is a smooth morphism
of affines with geometrically irreducible fibres of dimension $n$.

\medskip\noindent
The base case of the induction: $n = 0$.
As $X \to S$ is smooth with geometrically
irreducible fibres of dimension $0$ we see that $X \to S$ is an open
immersion, see
Descent, Lemma \ref{descent-lemma-universally-injective-etale-open-immersion}.
As $S$ is local and the closed point is in the image of $X \to S$
we conclude that $X = S$. Thus we see that $\mathcal{F}$ corresponds
to a finite flat $\mathcal{O}_{S, s}$ module. In this case the result
follows from
Algebra, Lemma \ref{algebra-lemma-finite-flat-local}
which tells us that $\mathcal{F}$ is in fact finite free.

\medskip\noindent
The induction step. Assume the result holds whenever the dimension
of the support in the closed fibre is $< n$. Write $S = \Spec(A)$,
$X = \Spec(B)$ and $\mathcal{F} = \widetilde{N}$ for some $B$-module
$N$. Note that $A$ is a local ring; denote its maximal ideal $\mathfrak m$.
Then $\mathfrak p = \mathfrak mB$ is the unique minimal prime lying over
$\mathfrak m$ as $X \to S$ has geometrically irreducible fibres. Finally,
let $\mathfrak q \subset B$ be the prime corresponding to $x$. By
Lemma \ref{lemma-induction-step}
we can choose a map
$$
\alpha : B^{\oplus r} \to N
$$
such that $\kappa(\mathfrak p)^{\oplus r} \to N \otimes_B \kappa(\mathfrak p)$
is an isomorphism. Moreover, as $N_{\mathfrak q}$ is $A$-flat the lemma
also shows that $\alpha$ is injective and that
$\Coker(\alpha)_{\mathfrak q}$ is $A$-flat.
Set $Q = \Coker(\alpha)$. Note that the support of $Q/\mathfrak mQ$
does not contain $\mathfrak p$. Hence it is certainly the case that
$\dim_{\mathfrak q}(\text{Supp}(Q/\mathfrak mQ)) < n$.
Combining everything we know about $Q$ we see
that the induction hypothesis applies to $Q$. It follows that there exists
an elementary \'etale morphism $(S', s) \to (S, s)$ such that the conclusion
holds for $Q \otimes_A A'$ over $B \otimes_A A'$ where
$A' = \mathcal{O}_{S', s'}$. After replacing $A$ by $A'$ we have an
exact sequence
$$
0 \to B^{\oplus r} \to N \to Q \to 0
$$
(here we use that $\alpha$ is injective as mentioned above)
of finite $B$-modules and we also get an element
$g \in B$, $g \not \in \mathfrak q$ such that
$Q_g$ is finitely presented over $B_g$ and flat over $A$. Since localization
is exact we see that
$$
0 \to B_g^{\oplus r} \to N_g \to Q_g \to 0
$$
is still exact. As $B_g$ and $Q_g$ are flat over $A$ we conclude that
$N_g$ is flat over $A$, see
Algebra, Lemma \ref{algebra-lemma-flat-ses},
and as $B_g$ and $Q_g$ are finitely presented over $B_g$ the same holds
for $N_g$, see
Algebra, Lemma \ref{algebra-lemma-extension}.
\end{proof}

\end{document}
